\documentclass{resume} % Use the custom resume.cls style

\usepackage[left=0.4in,top=0.4in,right=0.5in,bottom=0.4in]{geometry} % Document margins1
%%% \usepackage[utf8x]{inputenc}
\usepackage{epigrafica}
\usepackage[LGR,OT1]{fontenc}
\usepackage{hyperref}

\name{Jose Eduardo Ruiz Rosero} % Your name

\address{Sherbrooke, Canada} % Your address
% \address{Cali, Colombia} % Your address
% \address{Rio de Janeiro, Brazil} % Your address

\address{\href{https://www.linkedin.com/in/joseruiz19}{LinkedIn}, \href{https://scholar.google.com/citations?user=nE5XCmMAAAAJ&hl=en}{Google Scholar}, \href{https://github.com/joseruiz1989}{GitHub}}

\address{joseruiz1989@hotmail.com, \scriptsize{CV version: \today}} % Your phone number and email

\begin{document}
    \begin{rSection}{}
        \begin{center}
            MACHINE LEARNING ENGINEER - DEVELOPER - DATA SCIENTIST\\ \vspace*{0.3cm}
            {\normalfont 
            {\scriptsize I'm a Ph.D. with over 14 years of experience in AI, nanotech, web/software dev, and semiconductors. My skills include ML, Python, agile methods, containers, and collaboration across teams from diverse backgrounds, countries, and cultures for innovative solutions, driven by a passion for scientific progress.}
            }
        \end{center}
    \end{rSection}
    
    \begin{rSection}{}
    \end{rSection}
    
    %----------------------------------------------------------------------------------------
    %    CORE COMPETENCES
    %----------------------------------------------------------------------------------------
    \begin{rSection}{Core Competences}
        \begin{rSubsection}{}{}{}{}
            \item {\normalfont Over nine years of experience in Python programming language, especially with libraries/frameworks for machine learning, data processing, data visualization, image processing, multiprocessing. (Keras, TensorFlow, TensorRT, Detectron2, Mask-RCNN, Matplotlib, Plotly, Seaborn, PyTorch, OpenCV, Pillow, NumPy, Pandas, SciKit-Learn, Multiprocessing, Django).}
            \item {\normalfont High experience in heterogeneous computing with GPU, and Multi-core for simulation and machine learning, including clusters based on the Slurm system.}
            \item {\normalfont Significant experience with project management (PM) methodologies (Scrum), and PM tools (Jira, Slack, Trello).}
            \item {\normalfont Extensive experience working in multidisciplinary and multicultural groups, both in academics and industry.}
            \item {\normalfont Excellent technical documentation skills for the final customer and research publishing, using tools like MS Office suite, \LaTeX, Corel (Draw, Photo-Paint), Adobe (Premiere, Illustrator, Photoshop), OriginPro.}
            \item {\normalfont Over nine years of experience in researching semiconductor materials for solar cells, laser and photodetectors devices for academic research and industry.}
            \item {\normalfont Broad ability to understand, debug, and improve the existing source code.}
        \end{rSubsection}
    \end{rSection}
    
    %----------------------------------------------------------------------------------------
    %    EDUCATION SECTION
    %----------------------------------------------------------------------------------------
    \begin{rSection}{Education}

        {\bf Ph.D., Electrical Engineering} \hfill {\em 2015-2019} \\ 
        {\normalfont Pontifical Catholic University of Rio de Janeiro, PUC-Rio, Rio De Janeiro, Brazil}\\
        \underline{Thesis:} {\normalfont Precursors evaluation for GaInNAs growth by MOVPE technique for solar cells production.}\\ 
        Advisors: Ph.D. Patricia Lustoza de Souza (PUC-Rio).
        
        {\bf Ph.D. internship} \hfill {\em 2016-2018} \\ 
        {\normalfont Fraunhofer Institute for Solar Energy Systems, ISE-Fraunhofer, Freiburg im Breisgau, Germany}\\
        Advisors: Ph.D. Frank Dimroth (ISE-Fraunhofer).
        
        {\bf M.Sc., Electrical Engineering} \hfill {\em 2013-2015} \\ 
        {\normalfont Pontifical Catholic University of Rio de Janeiro, PUC-Rio, Rio De Janeiro, Brazil}\\
        \underline{Thesis:} {\normalfont Optical and morphological characterization of InAs quantum dots.}\\
        Advisors: Ph.D. Patricia Lustoza de Souza (PUC-Rio) and Ph.D. Mauricio Pamplona Pires (UFRJ).
                
        {\bf B.Sc., Electronic Engineer} \hfill {\em 2006-2013} \\ 
        {\normalfont University of Nariño, UDENAR, San Juan de Pasto, Colombia.}\\
        \underline{Thesis:} {\normalfont Characterization of Piezoelectric Disk Displacement as a Function of Voltage with Hysteresis Correction and Drift Compensation using Closed-Loop Control, employing Interferometry.} \\
        Advisor: M.Sc. Dario Fernando Fajardo (UDENAR).
    \end{rSection}
    
    \newpage
    
    \begin{rSection}{Postdoctoral Researcher}

        {\bf Postdoctoral Researcher} \hfill {\em 2024 - Present} \\ 
        {\normalfont University of Sherbrooke - Faculté de génie - Département de génie Électrique et Génie Informatique, Sherbrooke, Canada}\\
        \underline{Research:} {\normalfont Simulations and optimization of laser performance.} \\
        Advisors: Ph.D. Gwenaëlle Hamon (U.Sherbrooke).
        \begin{rSubsection}{}{}{}{}
            {\normalfont 
                \item Simulation and optimization of semiconductor laser structures, including material selection, layer thicknesses, and doping profiles in both n- and p-regions (waveguide and cladding), aiming to minimize operating voltage while maintaining device output power.
                \item Simulation and optimization of tunnel junctions (tunnel diodes) for multi-junction laser architectures.
                \item Conducting laser performance simulations, as well as participating in fabrication and characterization processes in a cleanroom environment.
                \item Mentoring graduate-level students, providing guidance and support throughout their research activities.
                \item Performing project management tasks, contributing to the planning and coordination of research projects.
            }
        \end{rSubsection}


        
        
        {\bf Postdoctoral Researcher} \hfill {\em 2021 - 2024} \\ 
        {\normalfont Federal University of Rio de Janeiro - Physics Institute, IF-UFRJ, Rio De Janeiro, Brazil}\\
        \underline{Research:} {\normalfont Nanostructured photodetectors development with artificial intelligence.} \\
        Advisors: Ph.D. Mauricio Pamplona Pires (UFRJ).
        \begin{rSubsection}{}{}{}{}
            {\normalfont 
                \item Created Python code for semiconductor structure simulations, focusing on photovoltaic photodetectors, to optimize their properties and responses.
                \item Utilized multiprocessing techniques to parallelize simulations on computers with multi-thread processors and on cloud computing platforms.
                \item Applied artificial intelligence, specifically evolutionary algorithms, to optimize simulations and propose novel device structures with improved performance.
                \item Developed extensive databases of simulation results to analyze and map the properties of the optimized devices for better understanding and insights.
                \item Assisted in preparing and organizing congresses, seminars, and scientific papers to disseminate research findings.
                \item Provided guidance and support to undergraduate, master's, doctoral students, and postdoctoral colleagues in their research projects.
                \item Translated Fortran code to Python, facilitating the implementation of complex physical process simulations in a user-friendly environment.
                \item Developed e-mulate software with a graphical user interface for electronic state calculations in quantum well heterostructures.
                \item Employed numerical methods like transfer matrix and Numerov algorithms to solve the Schrodinger equation for electronic states calculation.
                \item Validated e-mulate software by comparing calculated absorption spectra with experimental results for quantum Bragg mirror detectors, demonstrating its accuracy and reliability.
            }
        \end{rSubsection}

    \end{rSection}
    
    %----------------------------------------------------------------------------------------
    %    Awards / Honors
    %----------------------------------------------------------------------------------------

    \begin{rSection}{Awards / Honors}
        \begin{tabular}{ @{} >{\bfseries}l @{\hspace{6ex}} l }
            Ph.D. Scholarship & {\normalfont Brazilian National Council for Scientific and Technological Development} \\& {\normalfont (CNPq). 2015.}\vspace{0.5em} \\
            Master Scholarship & {\normalfont Brazilian Coordination for the Improvement of Higher Education} \\& {\normalfont Personnel (CAPES). 2013.}\vspace{0.5em} \\
            Funding for student research & {\normalfont Vicerectorate for Research, Postgraduate and International Relations,} \\& {\normalfont University of Nariño (VIPRI - UDENAR). "Alberto Quijano } \\& {\normalfont  Guerrero" Student Research Award, 2012.}\vspace{0.5em} \\
        \end{tabular}
    \end{rSection}

    \newpage
    %----------------------------------------------------------------------------------------
    %    Languages
    %----------------------------------------------------------------------------------------

    \begin{rSection}{Languages}
        \begin{tabular}{ @{} >{\bfseries}l @{\hspace{6ex}} l }
            Spanish & {\normalfont Native} \vspace{0.5em} \\
            Portuguese & {\normalfont Fluent} \vspace{0.5em} \\
            English & {\normalfont Professional working proficiency} \vspace{0.5em} \\
            French & {\normalfont Basic} \vspace{0.5em} \\
            German & {\normalfont Beginner} \vspace{0.5em} \\
        \end{tabular}
    \end{rSection}

    % \newpage
    %----------------------------------------------------------------------------------------
    %    Software Skills
    %----------------------------------------------------------------------------------------

    \begin{rSection}{Software Skills}
        \begin{tabular}{ @{} >{\bfseries}l @{\hspace{6ex}} l }

            Programming languages & {\normalfont Python, JavaScript, Bash Script.}\vspace{0.5em} \\
            
            Machine Learning Techniques & {\normalfont Data Acquisition, Cleaning, Preprocessing, Visualization, Analysis, }\\& {\normalfont Transformation, Augmentation, Synthetic Generation, Quality Assurance.} \\& {\normalfont Model Selection, Training, Hyperparameter Tuning, Transfer Learning, } \\& {\normalfont Model Evaluation, Model Deployment, Anomaly Detection, }\\& {\normalfont GPU Acceleration, Multiprocessing, High-Performance Computing}\vspace{0.5em} \\
            
            Machine Learning Libraries & {\normalfont Keras, TensorFlow, TensorRT, Detectron2, Matplotlib, Plotly, Seaborn,} \\& {\normalfont PyTorch, OpenCV, Pillow, NumPy, Pandas, Multiprocessing, SciKit-Learn.}\vspace{0.5em} \\
            
            Protocols and APIs & {\normalfont JSON, XML, HDF5.}\vspace{0.5em} \\
            
            Numerical computing & {\normalfont Matlab.}\vspace{0.5em} \\
            
            Operating systems & {\normalfont Linux, Windows.}\vspace{0.5em}\\
            
            Documentation packages & {\normalfont Latex, MS Office.}\vspace{0.5em}\\
            
            Other tools & {\normalfont CMG, Docker, Git (GitHub, GitLab), OriginPro, Slurm, Singularity, }
            \\& {\normalfont Photon Design (Harold, HaroldXY PicWave)}
            \\& {\normalfont Corel (Draw, Photo-Paint), Adobe (Premiere, Illustrator, Photoshop)}\vspace{0.5em}\\
        \end{tabular}
    \end{rSection}
    
    %----------------------------------------------------------------------------------------
    %    WORK EXPERIENCE SECTION
    %----------------------------------------------------------------------------------------
    
    \begin{rSection}{Experience}

        \begin{rSubsection}{Applied Computational Intelligence Laboratory (ICA), PUC-Rio.}{May  2019 - Present}{Machine Learning Researcher / Developer / Data Scientist}{Rio De Janeiro, Brazil}
            {\normalfont 
            \item Analyzed and Structured Datasets: Orchestrated the analysis and organization of diverse datasets for AI model training, spanning oil and gas reservoir data, extensive natural language processing datasets, and image/video datasets to facilitate various processing tasks.
            \item Data Quality Assurance: Implemented data cleaning and preprocessing procedures, ensuring the data's integrity and preparedness for AI model development, including meticulous data preparation for image and video datasets.
            \item Synthetic Data Generation: Enlarged dataset volumes through the generation of synthetic data, producing a wide array of images and videos to fortify model robustness in image and video processing applications.
            \item Machine Learning Model Training: Proficiently trained machine learning models, specializing in deep learning, for a range of tasks including well path prediction, reservoir behavior forecasting, natural language processing, and image/video analysis.
            \item Personnel Training: Provided comprehensive training for junior and semi-senior staff to enhance their proficiency in working with datasets and machine learning models, fostering their growth and skill development.
            \item Cutting-Edge Technology Exploration: Conducted comprehensive research on state-of-the-art AI technologies, focusing on advancements in image and video processing techniques, to remain at the forefront of industry innovations.
            \item High-Performance Computing: Employed parallel computing, cloud computing, and distributed computing using clusters and SLURM systems to expedite the processing of large-scale image and video datasets, consequently accelerating AI model training.
            \item GPU Optimization: Leveraged multiple GPUs across nodes and even employed multiple nodes to maximize the efficiency of processing, training, and inference tasks.
            \item Containerization: Utilized Docker and Singularity containers to streamline data processing workflows, ensuring seamless deployment across different environments, enhancing reproducibility and efficiency.
            \item Synthetic Dataset Development: Engineered extensive synthetic datasets tailored for text, image, and video processing tasks, bolstering the generalization capabilities of AI models.
            \item Framework Implementation: Employed frameworks such as "Detectron2" to train models for a variety of image analysis tasks, including detection, classification, segmentation, and semantic segmentation.
            \item Python and Django Web Development for API-Connected Applications: Backend and Frontend Development, Development of Python-based back-end systems using Django framework, coupled with front-end interfaces helped by bootstrap.
            \item Research Leadership: Took the helm in leading and managing research initiatives for text, numeric data, image, and video processing projects, providing expertise and guidance to achieve project objectives.
            \item Cross-Functional Collaboration: Collaborated closely with cross-functional teams, integrating AI-driven solutions into oil and gas exploration and production processes, natural language processing, and image/video analysis applications while concurrently nurturing the development of junior and semi-senior employees to enhance their skillsets in working with datasets and machine learning models.
            \item Reservoir simulation with CMG software, extraction and parser of simulation results, HDF5 data creation from the extracted data, dataset structuration and training of transformer model for a proxy system to predict reservoir production.
            }
        \end{rSubsection}

        \begin{rSubsection}{Lucro Colombia}{January 2022 - Present}{Machine Learning Consultant}{Colombia}
            {\normalfont 
            \item Machine Learning Projects: Conducted numerous Machine Learning (ML) implementation and development projects for clients, collaborating effectively with cross-functional teams.
            \item Understanding Business Requirements: Ensured a comprehensive understanding of business requirements to tailor ML solutions accordingly.
            \item End-to-End Model Development: Undertook end-to-end development and deployment of ML models, which included data preprocessing, feature engineering, and rigorous model evaluation to optimize accuracy and performance.
            \item Synthetic Data Creation: Played a vital role in the development of synthetic data and the expansion of datasets through data augmentation techniques.
            \item Dataset Amplification: Successfully increased dataset size and diversity, ultimately enhancing the robustness and generalization of the ML models delivered.
            \item TensorFlow 2 Object Detection API: Creation and configuration of Docker, model training, freezing and optimization, model export to ONNX format, conversion to float16, deployment, conversion to TensorRT format.
            }
        \end{rSubsection}

        \begin{rSubsection}{Applied Computational Intelligence Laboratory (ICA), PUC-Rio.}{March  2013 - April 2019}{Software and Web Developer}{Rio De Janeiro, Brazil}
            {\normalfont 
            \item Developed Code (Python, C\#, and MATLAB): Created code in Python, C\#, and MATLAB for various projects, including artificial intelligence decision support systems, optimization using genetic algorithms, and simulations of solar cells.
            \item Utilized Genetic Algorithms for Optimization: Employed genetic algorithms for optimization purposes, optimizing well allocation in the artificial intelligence decision support system.
            \item Scientific Research: Engaged in scientific research activities, exploring cutting-edge advancements in various fields of artificial intelligence, ensuring relevance and innovation in project developments.
            \item Data Analysis: Conducted in-depth data analysis, extracting valuable insights to inform decision-making processes and project refinements.
            \item Data Collection and Preprocessing: Orchestrated data collection initiatives and meticulously preprocessed data to ensure its readiness for machine learning and artificial intelligence applications.
            \item Web Development (HTML, CSS, and JavaScript): Took charge of web development tasks, utilizing HTML, CSS, and JavaScript to craft interactive and user-friendly web pages, ensuring an engaging online experience for users.
            \item Video and Image Processing: Specialized in video and image processing techniques, enhancing visual elements and ensuring high-quality outputs for a variety of projects.
            \item State-of-the-Art Research: Vigorously studied the state of the art in diverse fields of artificial intelligence, keeping abreast of the latest advancements to integrate them into project development.
            \item Presentation Preparation: Prepared engaging and informative presentations to convey project findings, research outcomes, and innovations to various stakeholders, facilitating knowledge dissemination and collaboration.
            }
        \end{rSubsection}

        \begin{rSubsection}{Escalab Academy}{Oct 2021 - Feb 2022 // May 2023 - Sep 2023}{Python, Data Science and Machine Learning Tutor}{Chile}
            {\normalfont 
            \item Python Programming Tutorship: Providing tutoring services focused on Python programming, emphasizing data manipulation, analysis, and machine learning techniques.
            \item Guiding Students: Responsibilities include guiding students through the complexities of Python, equipping them with essential skills for success in the evolving fields of data science and analytics.
            \item Creating a Supportive Learning Environment: A commitment to nurturing a conducive learning environment that fosters exploration and critical thinking, enabling learners to effectively apply their knowledge to real-world scenarios.
            \item Inspiring Future Data Enthusiasts: The primary goal is to inspire the upcoming generation of data enthusiasts, arming them with the tools and knowledge necessary to excel in today's technology-driven landscape.
            }
        \end{rSubsection}


        \begin{rSubsection}{Semiconductor Laboratory (LabSem), PUC-Rio}{March  2015 - April 2019}{Ph.D. researcher and student}{Rio De Janeiro, Brazil}
            {\normalfont 
            \item Semiconductor Materials and Solar Cell Research: The primary focus of the work revolved around semiconductor materials and their applications in solar cells using the MetalOrganic Vapour-Phase Epitaxy (MOVPE) reactor.
            \item Cutting-Edge Research: Conducted advanced research and fabrication of semiconductor materials, including GaAs, GaInAs, GaNAs, and GaInNAs, with a particular emphasis on their growth and characterization for solar cell development.
            \item Knowledge Impartation: Delivered knowledge on the utilization of Atomic Force Microscopy (AFM) for material characterization. Employed artificial intelligence techniques in microscopy image processing, script programming for data analysis, and effective data visualization methods for handling extensive datasets.
            \item Methodology Development: Developed methodologies, fabricated and validated samples, and conducted comprehensive analyses using AFM, Photoluminescence (PL), Electrochemical Capacitance Voltage (ECV), and in-situ MOVPE experiments to gain vital insights into material properties.
            \item Community Engagement: Actively interacted with the scientific community by curating and visually presenting research data and authored academic papers to disseminate findings.
            \item Multidisciplinary Approach: The Ph.D. journey embraced a multidisciplinary approach, encompassing semiconductor growth, material characterization, data analysis, and advanced imaging techniques. This approach contributed to advancing the understanding and application of semiconductor materials in solar cell technology.
            }
        \end{rSubsection}

        \begin{rSubsection}{Fraunhofer Institute for Solar Energy Systems (ISE-Fraunhofer)}{October 2016 - October 2018}{Ph.D. student and assistant researcher}{Freiburg im Breisgau, Germany}
            {\normalfont 
            \item Conducted semiconductor material research and growth for solar cell applications using MOVPE reactor.
            \item Performed characterization using AFM, PL, ECV, and in-situ MOVPE techniques.
            \item Handled the conceptualization, methodology, sample fabrication, validation, formal analysis, investigation, visualization, data curation, and paper writing within the scope of the work.
            \item Explored GaInNAs growth via metalorganic vapor phase epitaxy (MOVPE) with a focus on reducing carbon background in the material.
            \item Investigated the influence of various precursors and growth parameters on crystal morphology, carbon background, and element incorporation in the semiconductor.
            \item Implemented measures to reduce carbon background, ultimately leading to the growth of solar cells for material evaluation in devices.
            }
        \end{rSubsection}
        
        \begin{rSubsection}{Semiconductor Laboratory (LabSem), PUC-Rio}{March  2013 - April 2015}{Master researcher and student}{Rio De Janeiro, Brazil}
            {\normalfont 
            \item Semiconductor Characterization: Responsible for characterizing semiconductor structures with quantum dots through optical (PL) and morphological (AFM) methods, monitoring epitaxial growth of new semiconductor samples, and processing and analyzing AFM-acquired images.
            \item Data Analysis and Reporting: Structured data, analyzed trends in experimental results, prepared reports and presentations for knowledge dissemination, and stayed informed by reading scientific papers.
            \item Knowledge Sharing and Support: Conducted seminars on scientific articles and supported team members with measurements and sample/device characterization.
            }
        \end{rSubsection}

        \begin{rSubsection}{Pontifical Catholic University of Rio de Janeiro (PUC-Rio)}{July  2015 - December 2015}{Teaching Lecturer}{Rio De Janeiro, Brazil}
            {\normalfont
            \item Covered fundamental physical principles and various solar cell technologies.
            \item Conducted practical laboratory sessions to provide hands-on experience with solar panels.
            \item Assisted students in characterizing the properties of solar panels.
            \item Aimed to equip students with foundational knowledge and technical skills essential for working with photovoltaic solar energy.
            }
        \end{rSubsection}

        \begin{rSubsection}{Instrumentation and intelligent systems group, UDENAR}{October  2011 - February 2013}{Researcher}{San Juan de Pasto, Colombia}
            {\normalfont 
            \item Investigation into Michelson Interferometer: Conducted a comprehensive investigation into the principles of operation of a Michelson interferometer.
            \item Michelson Interferometer Setup Fabrication: Fabricated a Michelson interferometer setup based on acquired knowledge.
            \item Closed Loop Controller Design: Designed and implemented a closed-loop controller to precisely control the displacement of a piezoelectric disk at the nanometric scale, achieving a resolution of +- 1nm.
            }
        \end{rSubsection}
    \end{rSection}
    
    %----------------------------------------------------------------------------------------
    %    PUBLICATIONS
    %----------------------------------------------------------------------------------------
    

    \begin{rSection}{Publications}

        \begin{rSubsection}{Journals articles and Proceedings}{}{}{}
            \item X-Ray Diffraction of Infrared Quantum Bragg Mirror Detector Structures Designed by Genetic Algorithm \\ {\normalfont Jailson Oliveira Santana; José Eduardo Ruiz Rosero; Pedro Henrique Pereira; Luciana Dornelas Pinto; Rudy Massami Sakamoto Kawabata; Jéssica de Oliveira Lorenzi; Guilherme Monteiro Torelly; Patrícia Lustoza de Souza; Maurício Pamplona Pires; Germano Maioli Penello} In:  {\normalfont Journal of Integrated Circuits and Systems2026} \\ DOI: {\normalfont  \url{https://doi.org/10.29292/jics.v21i1.1198}}
            \item Genetic algorithm optimization of quantum Bragg mirror detectors \\ {\normalfont Jose E. Ruiz; Pedro H. Pereira; Guilherme M. Torelly; Patricia L. Souza; Mauricio P. Pires; Germano M. Penello} In: {\normalfont APL Electronic Devices. 2026} \\ DOI: {\normalfont  \url{https://doi.org/10.1063/5.0315540}}
            \item Mapping and Optimization of Oscillator Strength in Quantum Bragg Mirror Detectors as a Function of their Dimensions \\ {\normalfont Jose E. Ruiz; Pedro H. Pereira; Germano M. Penello; Guilherme M. Torelly; Patricia L. Souza; Mauricio P. Pires} In:  {\normalfont 2023 37th Symposium on Microelectronics Technology and Devices (SBMicro)} \\ DOI: {\normalfont  \url{https://doi.org/10.1109/SBMicro60499.2023.10302557}}
            \item Cognitive Search: a free information retrieval web service to coronavirus scientific papers \\ {\normalfont Cristian Enrique Munoz Villalobos; Leonardo Alfredo Forero Mendoza; Renato Sayão da Rocha; Jose Eduardo Ruiz; Harold De Mello; Marco Aurélio Cavalcanti Pacheco} In:  {\normalfont IEEE ColCACI 2023: IEEE Colombian Conference on Applications of Computational Intelligence. 2023} \\ DOI: {\normalfont  \url{https://doi.org/10.1109/ColCACI59285.2023.10225924}}
            \item A Free Web Service for Fast COVID-19 Classification of Chest X-Ray Images with Artificial Intelligence \\ {\normalfont Jose David Bermudez Castro; Jose E. Ruiz; Pedro Achanccaray Diaz; Smith Arauco Canchumuni; Cristian Muñoz Villalobos; Felipe Borges Coelho; Leonardo Forero Mendoza; Marco Aurelio C. Pacheco } In:  {\normalfont Computational Science and Its Applications - ICCSA 2022. ICCSA 2022. Lecture Notes in Computer Science, vol 13375. Springer, Cham. 2022} \\ DOI: {\normalfont  \url{https://doi.org/10.1007/978-3-031-10522-7_29}}
            \item e-mulate: a user-friendly software to calculate optoelectronic properties of quantum well systems. \\ {\normalfont Pereira, P. H.; Torelly, G. M.; Lacerda, M. S.; Souza, D. R.; Ruiz, J. E.; Souza, V. R.; Chipana, L. A.; Souza, P. L.; Penello, G. M.; Pires, M. P.} In:  {\normalfont 2021 35th Symposium on Microelectronics Technology and Devices (SBMicro), Campinas, Brazil. p. 1-4. 2021} \\ DOI: {\normalfont  \url{https://doi.org/10.1109/SBMicro50945.2021.9585740               }}
            \item Nitrogen and carbon incorporation in GaNxAs1-x grown in a showerhead MOVPE reactor. \\ {\normalfont Ruiz, J. E.; Lackner, D.; Souza, P.L.; Dimroth, F.; Ohlmann, J.} In:  {\normalfont Journal of Crystal Growth, v 557, 125998. 2021} \\ DOI: {\normalfont  \url{https://doi.org/10.1016/j.jcrysgro.2020.125998}}
            \item A free web service for fast COVID-19 classification of chest X-Ray images \\ {\normalfont Jose David Bermudez Castro; Ricardo Rei; Jose E. Ruiz; Pedro Achanccaray Diaz; Smith Arauco Canchumuni; Cristian Muñoz Villalobos; Felipe Borges Coelho; Leonardo Forero Mendoza; Marco Aurelio C. Pacheco} In:  {\normalfont arXiv 2009.01657. 2020} \\ DOI: {\normalfont  \url{https://doi.org/10.48550/arXiv.2009.01657}}
            \item Otimização da malha metálica de dedos coletores para o contato elétrico frontal de células solares por algoritmos genéticos \\ {\normalfont PEJENDINO, R. C.; MICHA, D. N.; RUIZ, J. E.; WEINER, E. C.; PIRES, M. P.; SOUZA, P. L.} In:  {\normalfont VII Congresso Brasileiro de Energia Solar, 2018, Gramado/RS. VII Congresso Brasileiro de Energia Solar - CBENS. 2018} \\ Available at: {\normalfont  \url{https://anaiscbens.emnuvens.com.br/cbens/article/view/259}}
            \item Optimization of contact grids for solar cells with genetic algorithms. \\ {\normalfont PEJENDINO, R. C.; RUIZ, J. E.; WEINER, E.; MICHA, D. N.; PIRES, M. P.; SOUZA, P. L.} In:  {\normalfont 32nd Symposium on Microelectronics Technology and Devices (SBMicro), Fortaleza, Brazil. p. 1. 2017} \\ DOI: {\normalfont  \url{https://doi.org/10.1109/SBMicro.2017.8112975}}
            \item Optimization of digital image processing to determine quantum dots height and density from atomic force microscopy. \\ {\normalfont RUIZ, J. E.; PACIORNIK, S.; PINTO, L. D.; PTAK, F.; PIRES, M. P.; SOUZA, P. L.} In:  {\normalfont ULTRAMICROSCOPY, v. 184, p. 234-241. 2017} \\ DOI: {\normalfont  \url{https://doi.org/10.1016/j.ultramic.2017.09.004}}
            \item Optimization in digital image processing to determine quantum dots heights from atomic force microscopy. \\ {\normalfont RUIZ, J. E.; SOUZA, P. L.; DORNELAS, L. P.; PIRES, M. P.} In:  {\normalfont XIV Brazil MRS Meeting, Rio de Janeiro. Proceedings of the XIV Brazil MRS Meeting, p. 222. 2015} \\ Available at: {\normalfont  \url{http://sbpmat.org.br/14encontro/anais/pdfs/plenary/ASEA.pdf}}
            \item Early nucleation stages of low density InAs quantum dots nucleation on GaAs by MOVPE. \\ {\normalfont TORELLY, G.; JAKOMIN, R.; PINTO, L. D.; PIRES, M. P.; RUIZ, J.; CALDAS, P. G.; PRIOLI, R.; XIE, H.; PONCE, F. A.; SOUZA, P. L.} In:  {\normalfont Journal of Crystal Growth, v. 434, p. 47-54. 2015} \\ DOI: {\normalfont  \url{https://doi.org/10.1016/j.jcrysgro.2015.10.031}}
        \end{rSubsection}

        \begin{rSubsection}{Conferences publications}{}{}{}
            \item Cybersecurity as a Pillar of Sustainable Development and Commercial Stability: Perspective from Digital Interdependence \\ {\normalfont Jose Eduardo Ruiz Rosero} In: {\normalfont XI UNADE International Symposium 2025: ``National Security and Trade: Strategy for Sustainable Development``, Dominican Republic. 2025} \\ Web event: {\normalfont \url{https://unade.edu.do/simposio-internacional/}} \\ YouTube live stream: {\normalfont \url{https://www.youtube.com/live/hy1DNI7qKws?si=EnFZyJxDX8pqYDld&t=23600}} % 2025
            \item X-Ray Diffraction of Infrared Quantum Bragg Mirror Detectors Structures Designed by Artificial Intelligence \\ {\normalfont Jailson Santana; Jose E. R. Rosero; Pedro Henrique Pereira; Luciana D. Pinto; Rudy M. S. Kawabata; Guilherme M. Torelly; Germano M. Penello; Patrícia L. Souza; Mauricio P. Pires} In:  {\normalfont SBMicro - 39th symposium on microelectronics technology and devices - Chip in the Jungle. 2025} \\ SBMicroProgram: {\normalfont  \url{https://sites.google.com/uea.edu.br/chip-in-the-jungle-2025/Program/program/sbmicro-program}} % 2025
            \item Optimizing Quantum Bragg Mirror Detector: A Genetic Algorithm Approach \\ {\normalfont Pedro H. Pereira; Jose E. Ruiz; Guilherme M. Torelly; Mauricio P. Pires; Patricia L. Souza; Germano M. Penello} In:  {\normalfont Mid-IR Optoelectronics: Materials and Devices XVII2025, Vienna, Austria. 2025} \\ web event: {\normalfont  \url{https://miomd2025.conf.tuwien.ac.at/}} \\ Conference Program and Publications: {\normalfont  \url{https://miomd2025.conf.tuwien.ac.at/preliminary-program/}} % 2025
            \item Ethical and Social Challenges in the Age of Human Enhancement: How Converging Technologies Could Transform Our Individuality and Behavior \\ {\normalfont Jose Eduardo Ruiz Rosero} In: {\normalfont XX International Humanities Congress ``Cyborg: The Human Condition in a Transhumanist World``, Bucaramanga, Colombia. 2023} \\ Web event: {\normalfont \url{https://humanidades.ustabuca.edu.co/index.php/extencion/congreso-2023}} \\ YouTube live stream: {\normalfont \url{https://www.youtube.com/live/KUPJM2GgcKo?feature=share&t=3369}} % 2023
            \item Panel Discussion \\ {\normalfont Jose Eduardo Ruiz Rosero} In: {\normalfont XX International Humanities Congress ``Cyborg: The Human Condition in a Transhumanist World``, Bucaramanga, Colombia. 2023} \\ Web event: {\normalfont \url{https://humanidades.ustabuca.edu.co/index.php/extencion/congreso-2023}} \\ YouTube live stream: {\normalfont \url{https://www.youtube.com/live/pTePCup8ukg?feature=share&t=4014}} % 2023
            \item Quantum Bragg Mirror Detector Optimization Applying Genetic Algorithms \\ {\normalfont Jose Ruiz; Pedro H. Pereira; Germano M. Penello; Guilherme M. Torelly; Patricia L. Souza; Mauricio P. Pires} In:  {\normalfont Autumn Meeting of the Brazilian Physical Society, Ouro Preto. 2023} \\ Available at: {\normalfont  \url{https://sec.sbfisica.org.br/eventos/eosbf/2023/sys/resumos/R0784-1.pdf}} % 2023
            \item Mapping and Optimization of Oscillator Strength in Quantum Bragg Mirror Detectors as a Function of their Dimensions \\ {\normalfont Jose Ruiz; Pedro H. Pereira; Germano M. Penello; Guilherme M. Torelly; Patricia L. Souza; Mauricio P. Pires} In:  {\normalfont SBMicro - 37th symposium on microelectronics technology and devices - Chip in Rio. 2023} % 2023
            \item Dilute Nitride GaInNAs Growth with Low Decomposition Temperature Precursors by MOVPE. \\ {\normalfont RUIZ, J. E.; OHLMANN, J.; DIMROTH, F.; LACKNER, D.} In:  {\normalfont DGKK “Deutsche Gesellschaft für Kristallwachstum und Kristallzüchtung” - Freiburg im Breisgau - Germany. 2017} \\ Available at: {\normalfont  \url{https://www.ise.fraunhofer.de/content/dam/ise/de/downloads/pdf/DGKK_2017_Homepage-Programm.pdf}} % 2017
            \item Simulação e otimização de células solares simples e multijunção com algoritmos genéticos. \\ {\normalfont PACHECO, M. A. C.; RUIZ, J. E.; GALVÃO, P.} In:  {\normalfont 2017. In: XIX CECEMM: Congresso dos Estudantes de Ciência e Engenharia de Materiais do MERCOSUL. 2017} % 2017
            \item Dependence of the strain of InAs/InGaAlAs QDs on their size. \\ {\normalfont RUIZ, J. E.; PIRES, M. P.; SOUZA, P. L.} In:  {\normalfont XIV Brazil MRS Meeting, 2015, Rio de Janeiro. Proceedings of the XIV Brazil MRS Meeting.2015} \\ Available at: {\normalfont  \url{http://sbpmat.org.br/14encontro/anais/pdfs/plenary/ASEV.pdf}} % 2015
            \item From InAs extended monolayer flat 2D terraces to 3D islands grown on GaAs substrates. \\ {\normalfont Torelly, G; RUIZ, J.; JAKOMIN, R.; DORNELAS, L. P.; PIRES, M. P.; CALDAS, P. G.; PRIOLI, R.; XIE, H.; PONCE, F. A.; SOUZA, P. L.} In:  {\normalfont In: XIV Brazil MRS Meeting, Rio de Janeiro.2015} % 2015
        \end{rSubsection}

        \begin{rSubsection}{Chapters in books}{}{}{}
            \item Otimização da malha metálica de dedos coletores para o contato elétrico frontal de células solares por algoritmos genéticos. \\ {\normalfont Roberto Carlos Pejendino Jojoa; Daniel Neves Micha; Jose Eduardo Ruiz Rosero; Eleonora Cominato Weiner; Mauricio Pamplona Pires; Patrícia Lustoza de Souza} In:  {\normalfont Energia no Brasil. 1ed.: Editora Poisson, v. 1, p. 89-101. 2019} \\ DOI: {\normalfont  \url{https://doi.org/10.5935/978-85-7042-107-4        }}
        \end{rSubsection}

        \begin{rSubsection}{Patents}{}{}{}
            \item Method for Extracting and Structuring Information \\ {\normalfont } In:  {\normalfont BR 10 2021 023977-8. Issued Nov 26, 2021 - Journal number: 2662 - Section VI patents. 2021} \\ Available at: {\normalfont  \url{http://revistas.inpi.gov.br/rpi/}}
        \end{rSubsection}
        
    \end{rSection}
    
\end{document}
